\hypertarget{abstract}{%
\section*{Abstract}\addcontentsline{toc}{section}{Abstract}\label{abstract}}

A finite arithmetic selector can be transported into a rational signal
projector and a complex projective measurement by one explicitly
normalized Fourier transformation. This paper gives the full proof chain
and follows a single exact twelve-dimensional example through every
representation. The arithmetic input is the coprime residue set modulo
twelve. Its Fourier synthesis is the integer-valued Ramanujan signal.
The associated circulant matrix is an orthogonal projector of rank four.
A normalized version of the signal is a quantum state, and an inverse
Fourier transformation recovers the arithmetic support. An unsigned
state has the same probabilities in the original basis yet zero
probability on the selected support after transformation. The two states
are orthogonal. A continuous relative-phase family makes that difference
visible without changing the initial probability histogram.

The general finite Fourier, Möbius, divisor-projector, and polyphase
identities are proved in the text. The formal companion covers the
concrete modulus-twelve realization, including the sound interpretation
of exact rational coordinates in the ordinary complex numbers, both
unitarity identities, projector transport, state preparation and
recovery, and the phase-sensitive probability conclusions. The later
inherited checkpoint contains ninety-six named theorems across five
production modules and preserves the twenty-two original bridge
statements. The current edition separates that archived compiler and
kernel evidence from current source audits and Python replay. A fresh
Lean build was not available for this edition.

Ramanujan sums, Fourier orthogonality, rational cyclic projectors, and
Ramanujan signal subspaces are classical. The contribution of this
presentation is an exact, inspectable realization chain, its
theorem-to-code concordance, and its explicit account of the information
retained or discarded at each map. It makes no circuit-efficiency or
quantum-advantage claim.

\hypertarget{the-mathematical-object-being-transported}{%
\section{The mathematical object being
transported}\label{the-mathematical-object-being-transported}}

The word \emph{bridge} is useful only when its arrows have definitions.
In the present example the common object is a selector: a diagonal
operator that retains certain Fourier coordinates and discards the
others. Arithmetic supplies the selected labels. Fourier synthesis
expresses that selector in sample coordinates. Complex normalization
then supplies a unitary transformation and the state vectors used in a
quantum interpretation.

The basic identity is

\[
P=FDF^\dagger.
\]

Here \(D\) is a diagonal zero-or-one matrix, \(F\) is unitary, and
\(\dagger\) denotes conjugate transpose. Consequently \(P\) is an
orthogonal projector. The transformation \(F\) is reversible. The
selection \(P\) need not be. A complete decomposition consists of both
\(Px\) and \((I-P)x\), whose sum recovers \(x\).

This distinction determines what can be transported. A Fourier basis
change preserves dimension, inner products, norms, and the entire
vector. A projector preserves only its selected component. Measurement
probabilities discard still more information. The same exact operator
can appear in arithmetic, signal processing, and finite-dimensional
quantum theory, while those uses retain different observables and have
different implementation costs.

We work over \(\mathbb C\) for Fourier transformations and use
\(\mathbb Q\) for the rational kernels. All index sets in the main
finite construction are specified, and no limit argument is needed. A
packet in this paper means a component together with its stated
realization map. This definition does not turn every kind of component
into the same object. Character-order packets and coefficient-residue
lanes will be distinguished explicitly in Section 6.

\hypertarget{claim-labels-and-evidence-classes}{%
\subsection{Claim labels and evidence
classes}\label{claim-labels-and-evidence-classes}}

The labels QA01 through QA20 refer to mathematical statements in this
paper. They are not a count of Lean declarations. A single claim can use
several formal lemmas, and a general theorem can have a formally checked
specialization rather than a generic Lean proof.

We use three evidence descriptions. A \emph{general prose theorem} has a
proof below for all inputs satisfying its hypotheses. A \emph{formal
specialization} is an actual declaration in the preserved modulus-twelve
source. A \emph{finite computational check} is a bounded verification of
an implementation. These descriptions can support the same identity, but
they are not interchangeable. The Lean guide gives the exact declaration
mapping.

The central example uses the vectors

\[
\begin{aligned}
c&=(4,0,2,0,-2,0,-4,0,-2,0,2,0),\\
b&=(5,2,-1,2,5,2,-1,2,5,2,-1,2),\\
x&=(9,2,1,2,3,2,-5,2,3,2,1,2).
\end{aligned}
\]

These are the literal coefficients of the formal source, and \(x=c+b\).
Keeping one example fixed makes every arrow independently checkable.

\hypertarget{fourier-conventions-and-the-finite-proof-chain}{%
\section{Fourier conventions and the finite proof
chain}\label{fourier-conventions-and-the-finite-proof-chain}}

Let \(q\) be a positive integer. Write \(I_q=\{0,1,\ldots,q-1\}\),
interpreting indices modulo \(q\), and put

\[
\zeta_q=e^{2\pi i/q}.
\]

A signal is a vector \(x:I_q\to\mathbb C\). We write \(\operatorname{Id}_q\) for the identity matrix of dimension \(q\), and \(I\) for an identity matrix when its dimension is clear. Our inner product is
conjugate-linear in the first argument:

\[
\langle x,y\rangle=\sum_{n=0}^{q-1}\overline{x(n)}y(n),
\qquad \|x\|^2=\langle x,x\rangle.
\]

The unscaled negative-sign discrete Fourier transform is

\[
\widehat{x}(k)=\sum_{n=0}^{q-1}x(n)\zeta_q^{-kn}.
\]

The positive-sign normalized Fourier matrix is

\[
(F_q)_{na}=q^{-1/2}\zeta_q^{na}.
\]

Thus \(F_q\) synthesizes a signal from normalized frequency coordinates,
while \(F_q^\dagger x=\widehat{x}/\sqrt q\) analyzes it. This convention
explains all signs and all square-root factors below. The negative-sign
transform and the positive-sign synthesis are adjoints of one another.

\hypertarget{theorem-qa01.-the-order-of-a-finite-phase}{%
\subsection*{Theorem QA01. The order of a finite
phase}\label{theorem-qa01.-the-order-of-a-finite-phase}}

For \(a\in I_q\), the multiplicative order of \(\zeta_q^a\) is

\[
\frac q{\gcd(a,q)}.
\]

\textbf{Proof.} The equation \((\zeta_q^a)^t=1\) holds exactly when
\(q\) divides \(at\). Write \(g=\gcd(a,q)\), \(q=gq'\), and \(a=ga'\),
with \(\gcd(a',q')=1\). The divisibility condition is \(q'\mid a't\),
which is equivalent to \(q'\mid t\). Its smallest positive solution is
\(t=q'\). This includes \(a=0\): then \(g=q\) and the order is one.
\(\square\)

\hypertarget{theorem-qa02.-character-orthogonality}{%
\subsection*{Theorem QA02. Character
orthogonality}\label{theorem-qa02.-character-orthogonality}}

For any integer \(h\),

\[
\sum_{n=0}^{q-1}\zeta_q^{hn}
=\begin{cases}q,&q\mid h,\\0,&q\nmid h.\end{cases}
\]

\textbf{Proof.} If \(q\mid h\), every summand is one. Otherwise
\(r=\zeta_q^h\ne1\) and \(r^q=1\). Multiplying the sum by \(1-r\) gives
\(1-r^q=0\), so the sum is zero. This argument also covers \(q=1\),
where only the first case occurs. \(\square\)

\hypertarget{theorem-qa03.-unitarity-inversion-and-parseval}{%
\subsection*{Theorem QA03. Unitarity, inversion, and
Parseval}\label{theorem-qa03.-unitarity-inversion-and-parseval}}

For every positive \(q\),

\[
F_q^\dagger F_q=\operatorname{Id}_q,
\qquad F_qF_q^\dagger=\operatorname{Id}_q.
\]

Consequently,

\[
x(n)=\frac1q\sum_{k=0}^{q-1}\widehat{x}(k)\zeta_q^{kn},
\qquad
\sum_n|x(n)|^2=\frac1q\sum_k|\widehat{x}(k)|^2.
\]

\textbf{Proof.} The \((a,b)\) entry of the first matrix product is

\[
\frac1q\sum_{n=0}^{q-1}\zeta_q^{n(b-a)}.
\]

By QA02 it is one when \(a=b\) and zero otherwise. The \((n,m)\) entry
of the second product is the same orthogonality sum with indices
interchanged. Since \(F_q^\dagger x=\widehat{x}/\sqrt q\), the identity
\(F_qF_q^\dagger x=x\) is the displayed inverse formula. Finally, a
unitary matrix preserves inner products:

\[
\|F_q^\dagger x\|^2
=x^\dagger F_qF_q^\dagger x
=x^\dagger x.
\]

Replacing \(F_q^\dagger x\) by \(\widehat{x}/\sqrt q\) proves Parseval.
\(\square\)

These proofs use finite sums, ordinary complex multiplication, and the
positive real square root of \(q\). No asymptotic approximation enters.
At \(q=1\) the Fourier matrix is the one-by-one identity, and all
identities have their literal scalar meanings.

\hypertarget{coprimality-becomes-an-integer-signal}{%
\section{Coprimality becomes an integer
signal}\label{coprimality-becomes-an-integer-signal}}

Let

\[
U_q=\{a\in I_q:\gcd(a,q)=1\},
\qquad \varphi(q)=|U_q|.
\]

The Ramanujan sum is

\[
c_q(n)=\sum_{a\in U_q}\zeta_q^{an},
\qquad n\in\mathbb Z.
\]

It adds the primitive \(q\)th-root phases selected by QA01. Its
definition is complex, but its value is an integer. That integrality is
the step that makes the sample-coordinate projector rational. Ramanujan
sums and their divisor formulas are classical arithmetic constructions
\cite{Ramanujan1918}.

\hypertarget{theorem-qa04.-muxf6bius-formula-and-integrality}{%
\subsection*{Theorem QA04. Möbius formula and
integrality}\label{theorem-qa04.-muxf6bius-formula-and-integrality}}

Let \(\mu\) be the Möbius function, with \(\mu(1)=1\). For every
positive \(q\) and every integer \(n\),

\[
\boxed{c_q(n)=\sum_{d\mid\gcd(q,n)}d\,\mu(q/d)\in\mathbb Z.}
\]

In particular \(c_q\) has period \(q\), is real-valued, satisfies
\(c_q(-n)=c_q(n)\), and has \(c_q(0)=\varphi(q)\).

\textbf{Proof.} For a positive integer \(g\),

\[
\sum_{d\mid g}\mu(d)
=\begin{cases}1,&g=1,\\0,&g>1.\end{cases}
\]

To see this, only squarefree divisors contribute. If \(g\) has \(r\)
distinct prime factors, the sum is \((1-1)^r\), with value one when
\(r=0\). Therefore

\[
\mathbf1_{\gcd(a,q)=1}=\sum_{d\mid\gcd(a,q)}\mu(d).
\]

Substitute this into the phase sum over all \(a\in I_q\), exchange the
finite sums, and write \(a=db\) in the terms where \(d\mid a\):

\[
\begin{aligned}
c_q(n)
&=\sum_{d\mid q}\mu(d)
  \sum_{b=0}^{q/d-1}e^{2\pi ibn/(q/d)}\\
&=\sum_{\substack{d\mid q\\(q/d)\mid n}}\mu(d)\frac qd.
\end{aligned}
\]

The second equality is QA02. Replacing \(q/d\) by \(e\) gives the
displayed divisor sum. It is an integer. Periodicity follows already
from the phase definition. The divisor formula is unchanged by
\(n\mapsto -n\), so \(c_q(-n)=c_q(n)\). At \(n=0\), every selected phase
is one, giving \(c_q(0)=|U_q|\). \(\square\)

At \(q=1\), \(U_1=\{0\}\) because \(\gcd(0,1)=1\). The phase sum and the
divisor formula both give \(c_1(n)=1\). This boundary case is included
in the proofs even though the inherited arithmetic regression range
begins at modulus two.

\hypertarget{theorem-qa05.-exact-spectral-support-and-energy}{%
\subsection*{Theorem QA05. Exact spectral support and
energy}\label{theorem-qa05.-exact-spectral-support-and-energy}}

Over one period,

\[
\boxed{\widehat{c_q}(k)=q\,\mathbf1_{k\in U_q},\qquad k\in I_q.}
\]

Moreover,

\[
\sum_{n=0}^{q-1}|c_q(n)|^2=q\varphi(q).
\]

\textbf{Proof.} Substitute the phase definition and exchange two finite
sums:

\[
\widehat{c_q}(k)
=\sum_{a\in U_q}\sum_{n=0}^{q-1}\zeta_q^{(a-k)n}.
\]

QA02 makes the inner sum \(q\) exactly when \(a=k\), and zero otherwise.
The energy identity follows from QA03:

\[
\|c_q\|^2
=\frac1q\sum_k|\widehat{c_q}(k)|^2
=\frac1q\varphi(q)q^2.
\] \(\square\)

The selector has been recovered by an identity. A numerical spectrum can
illustrate the result, but neither a tolerance nor a peak-finding rule
defines the arithmetic support.

\hypertarget{rational-character-order-projectors}{%
\section{Rational character-order
projectors}\label{rational-character-order-projectors}}

It is helpful to distinguish the size of an ambient signal from the
order of a character. Fix a positive ambient length \(N\). For each
divisor \(d\) of \(N\), define

\[
S_d=\{a\in I_N:N/\gcd(a,N)=d\}.
\]

The set \(S_d\) contains the labels of characters of exact order \(d\).
Put \(D_d=\operatorname{diag}(\mathbf1_{a\in S_d})\) and

\[
P_d^{(N)}=F_ND_dF_N^\dagger.
\]

The superscript records the ambient dimension. In particular, the kernel
normalization below is \(1/N\), even when \(d<N\).

\hypertarget{theorem-qa06.-rational-kernels-and-complete-divisor-decomposition}{%
\subsection*{Theorem QA06. Rational kernels and complete divisor
decomposition}\label{theorem-qa06.-rational-kernels-and-complete-divisor-decomposition}}

For every \(d\mid N\),

\[
\boxed{(P_d^{(N)})_{nm}=\frac1N c_d(n-m).}
\]

The matrices have rational entries and satisfy

\[
\begin{aligned}
(P_d^{(N)})^\dagger&=P_d^{(N)},\\
P_d^{(N)}P_e^{(N)}&=\delta_{de}P_d^{(N)},\\
\sum_{d\mid N}P_d^{(N)}&=\operatorname{Id}_N,\\
\operatorname{rank}P_d^{(N)}&=\operatorname{tr}P_d^{(N)}=\varphi(d).
\end{aligned}
\]

\textbf{Proof.} A label in \(S_d\) has the unique form \(a=(N/d)b\),
where \(0\le b<d\) and \(\gcd(b,d)=1\). Indeed,

\[
\gcd((N/d)b,N)=(N/d)\gcd(b,d).
\]

Direct multiplication therefore gives

\[
\begin{aligned}
(F_ND_dF_N^\dagger)_{nm}
&=\frac1N\sum_{a\in S_d}\zeta_N^{a(n-m)}\\
&=\frac1N\sum_{b\in U_d}\zeta_d^{b(n-m)}
=\frac1N c_d(n-m).
\end{aligned}
\]

QA04 proves rationality. Each \(D_d\) is a self-adjoint diagonal
projector. The sets \(S_d\) partition \(I_N\), because every character
has exactly one order dividing \(N\). Hence \(D_dD_e=\delta_{de}D_d\)
and \(\sum_{d\mid N}D_d=\operatorname{Id}_N\). Conjugating these identities by the
unitary \(F_N\) proves the projector identities. The rank is unchanged
under multiplication by invertible matrices, so it equals the number of
selected diagonal entries, \(|S_d|=\varphi(d)\). The diagonal kernel
entries are all \(c_d(0)/N=\varphi(d)/N\), whose sum is \(\varphi(d)\).
\(\square\)

This orthogonal divisor-shell decomposition has established antecedents
in Ramanujan subspaces and filter banks \cite{Vaidyanathan2014FIR,Vaidyanathan2014I,Vaidyanathan2014II}. The general theorem above is
proved independently in this text. The formal companion certifies the
concrete \(N=d=12\) projector and its transport identity, rather than an
arbitrary-divisor Lean theorem.

\hypertarget{a-shift-operator-realization-and-its-sign}{%
\subsection{A shift-operator realization and its
sign}\label{a-shift-operator-realization-and-its-sign}}

Let \(A\) be the cyclic shift \((Ax)(n)=x(n-1)\). A Fourier column
\(f_a(n)=N^{-1/2}\zeta_N^{na}\) satisfies \(Af_a=\zeta_N^{-a}f_a\).
Therefore

\[
E_a=\frac1N\sum_{\ell=0}^{N-1}\zeta_N^{a\ell}A^\ell
\]

acts as the identity on \(f_a\) and as zero on every other column, by
QA02. Summing \(E_a\) over \(S_d\) gives

\[
P_d^{(N)}=\frac1N\sum_{\ell=0}^{N-1}c_d(\ell)A^\ell.
\]

This formula also follows entrywise because \((A^\ell)_{nm}\) is one
exactly when \(n-m=\ell\) modulo \(N\). A group-algebra convention with
the opposite shift sign writes individual idempotents with the opposite
phase sign. The divisor shells are unchanged under \(a\mapsto -a\), and
\(c_d(-\ell)=c_d(\ell)\), so the rational shell formula remains the
same.

Proposition 21.11 of \emph{Stratified Universality for Resolved
Representations}, v1.16.0, gives the rational cyclic order-projector
formula in the group algebra. Boundary 21.12 distinguishes it from
formal-series coefficient-residue selection \cite{KomorebiStratified16}.
The present shift calculation makes the representation and the sign
explicit.

\hypertarget{unions-of-order-classes}{%
\subsection{Unions of order
classes}\label{unions-of-order-classes}}

If \(\mathcal D\) is any set of divisors of \(N\), then

\[
P_{\mathcal D}=\sum_{d\in\mathcal D}P_d^{(N)}
\]

is a rational orthogonal projector of rank
\(\sum_{d\in\mathcal D}\varphi(d)\). Its complement is the sum over the
remaining divisors. Individual split Fourier projectors can have
nonrational entries. Summing complete order classes is a sufficient
mechanism for rationality. No converse classification of every rational
selector is needed for the present bridge.

\hypertarget{selection-residuals-and-information-loss}{%
\section{Selection, residuals, and information
loss}\label{selection-residuals-and-information-loss}}

For any set \(S\subseteq I_N\), let \(D_S\) be its diagonal indicator
and \(P_S=F_ND_SF_N^\dagger\). This construction does not require \(S\)
to be an order class. Rationality is special, but orthogonal projection
is general.

\hypertarget{theorem-qa07.-projector-geometry-and-exact-residual-reconstruction}{%
\subsection*{Theorem QA07. Projector geometry and exact residual
reconstruction}\label{theorem-qa07.-projector-geometry-and-exact-residual-reconstruction}}

For every \(x\in\mathbb C^N\),

\[
\begin{aligned}
x&=P_Sx+(I-P_S)x,\\
\langle P_Sx,(I-P_S)x\rangle&=0,\\
\|x\|^2&=\|P_Sx\|^2+\|(I-P_S)x\|^2.
\end{aligned}
\]

Furthermore, \(P_S\) is injective exactly when \(S=I_N\). If \(S\) is
proper, its kernel has dimension \(N-|S|\).

\textbf{Proof.} The first equality is distributivity. Since
\(P_S^\dagger=P_S\) and \(P_S^2=P_S\),

\[
\langle P_Sx,(I-P_S)x\rangle
=x^\dagger P_S(I-P_S)x=0.
\]

Expanding the norm square of the orthogonal sum proves the energy
identity. In Fourier coordinates, \(P_S\) is \(D_S\), whose kernel
consists exactly of the vectors supported outside \(S\). Its dimension
is \(N-|S|\). This kernel is zero exactly when every label is selected.
\(\square\)

This gives a precise reconstruction contract. The pair of packets
\((P_Sx,(I-P_S)x)\) retains the complete vector, subject to the
compatibility condition that the first lies in the range and the second
in the kernel. Keeping only the first packet is an exact projection but
an incomplete representation of an arbitrary input.

For character-order packets, QA06 gives the stronger complete
decomposition

\[
x=\sum_{d\mid N}P_d^{(N)}x,
\qquad
\|x\|^2=\sum_{d\mid N}\|P_d^{(N)}x\|^2.
\]

The decomposition is a change in organization rather than a reduction in
the dimension of the whole signal. Smaller stored selectors do not make
an arbitrary signal contain less information. Compression requires a
restriction on the admissible signals or an explicit approximation
criterion.

\hypertarget{coefficient-residue-lanes-and-polyphase-transport}{%
\section{Coefficient-residue lanes and polyphase
transport}\label{coefficient-residue-lanes-and-polyphase-transport}}

There is a second exact decomposition that is common in Komorebi
calculations and in signal processing. It separates \emph{positions} by
residue class rather than separating \emph{Fourier characters} by order.
Both decompositions use projectors, but they select different objects.

\hypertarget{formal-coefficient-lanes}{%
\subsection{Formal coefficient
lanes}\label{formal-coefficient-lanes}}

Let \(L\ge1\) and let \(f(z)=\sum_{n\ge0}a_nz^n\) be a formal power
series over a field containing a primitive \(L\)th root \(\omega\) and
\(1/L\). For \(0\le r<L\), define

\[
E_rf(t)=\sum_{m\ge0}a_{r+Lm}t^m,
\qquad
R_rf(z)=z^rE_rf(z^L).
\]

The first expression addresses the lane in a new variable. The second
places that lane back at its original coefficient positions.

\Needspace{13\baselineskip}
\hypertarget{theorem-qa08.-exact-coefficient-residue-decomposition}{%
\subsection*{Theorem QA08. Exact coefficient-residue
decomposition}\label{theorem-qa08.-exact-coefficient-residue-decomposition}}

The coefficient-residue maps satisfy

\[
\begin{aligned}
R_rf(z)&=\frac1L\sum_{j=0}^{L-1}\omega^{-rj}f(\omega^jz),\\
R_rR_s&=\delta_{rs}R_r,\\
\sum_{r=0}^{L-1}R_rf&=f.
\end{aligned}
\]

\textbf{Proof.} The coefficient of \(z^n\) in the first right-hand side
is

\[
\frac{a_n}{L}\sum_{j=0}^{L-1}\omega^{j(n-r)}.
\]

Character orthogonality makes it \(a_n\) when \(n\equiv r\pmod L\) and
zero otherwise. This is precisely the coefficient of \(R_rf\). Each
coefficient position belongs to a unique lane, so the projector and
completeness identities follow coefficientwise. The equality is formal,
with no analytic convergence hypothesis. \(\square\)

For an infinite analytic series, the same identity holds wherever all
displayed evaluations are defined, but the formal theorem does not need
that interpretation. The lane index is an exponent residue. It is not a
multiplicative-order class of a Fourier character.

\hypertarget{finite-polyphase-lanes}{%
\subsection{Finite polyphase lanes}\label{finite-polyphase-lanes}}

Now suppose \(N=LM\) with positive integers \(L\) and \(M\). For a
vector \(x\in\mathbb C^N\), define its \(L\) lanes by

\[
x_r(m)=x(r+Lm),
\qquad 0\le r<L,\quad 0\le m<M.
\]

The addressing map \(T:x\mapsto(x_0,\ldots,x_{L-1})\) is bijective.
Interleaving the lanes reconstructs \(x\), and

\[
\sum_{n=0}^{N-1}|x(n)|^2
=\sum_{r=0}^{L-1}\sum_{m=0}^{M-1}|x_r(m)|^2.
\]

This energy identity is an exact reindexing. It is useful to preserve it
separately from the spectral aliasing identity that follows.

\Needspace{19\baselineskip}
\hypertarget{theorem-qa09.-polyphase-fourier-identities-and-decimation-aliasing}{%
\subsection*{Theorem QA09. Polyphase Fourier identities and decimation
aliasing}\label{theorem-qa09.-polyphase-fourier-identities-and-decimation-aliasing}}

Write \(\widehat{x_r}\) for the unscaled \(M\)-point transform of lane
\(r\). Then for \(k\in I_N\),

\[
\boxed{\widehat{x}(k)
=\sum_{r=0}^{L-1}\zeta_N^{-kr}\widehat{x_r}(k\bmod M).}
\]

Conversely, for \(\ell\in I_M\),

\[
\boxed{\widehat{x_r}(\ell)
=\frac1L\sum_{j=0}^{L-1}
\zeta_N^{r(\ell+jM)}\widehat{x}(\ell+jM).}
\]

\textbf{Proof.} Partition the \(N\) samples by their unique
representation \(n=r+Lm\):

\[
\begin{aligned}
\widehat{x}(k)
&=\sum_{r=0}^{L-1}\sum_{m=0}^{M-1}
x_r(m)\zeta_N^{-k(r+Lm)}\\
&=\sum_r\zeta_N^{-kr}\sum_mx_r(m)\zeta_M^{-km}.
\end{aligned}
\]

The inner sum is \(\widehat{x_r}(k\bmod M)\), proving the first
identity. For the second, insert Fourier inversion into the lane
transform:

\[
\begin{aligned}
\widehat{x_r}(\ell)
&=\sum_{m=0}^{M-1}x(r+Lm)\zeta_M^{-\ell m}\\
&=\frac1N\sum_{k=0}^{N-1}\widehat{x}(k)\zeta_N^{kr}
\sum_{m=0}^{M-1}\zeta_M^{(k-\ell)m}.
\end{aligned}
\]

QA02 makes the inner sum \(M\) exactly when \(k=\ell+jM\) for some
\(0\le j<L\). Since \(M/N=1/L\), this is the second displayed identity.
\(\square\)

Every lane transform combines an alias class of \(L\) full-spectrum
coefficients. One decimated lane generally cannot distinguish those
coefficients. All \(L\) offsets together retain the original signal. For
a fixed alias class \(\ell\), the equations across \(r\) are an
invertible Fourier mixing system after the nonzero phases
\(\zeta_N^{r\ell}\) are removed. The apparent collapse is resolved by
preserving the complete collection of offsets.

\hypertarget{why-the-two-projectors-must-not-be-conflated}{%
\subsection{Why the two projectors must not be
conflated}\label{why-the-two-projectors-must-not-be-conflated}}

A coefficient-residue selector keeps samples such as
\(x(0),x(3),x(6),x(9)\). A character-order projector keeps spectral
modes such as \(1,5,7,11\). The former is diagonal in sample
coordinates; the latter is diagonal in frequency coordinates. Under a
Fourier basis change, a sample-residue mask generally produces a
different operator, with its own alias structure.

A valid comparison therefore specifies the conjugating matrix and the
chosen selector. The fact that both constructions use roots of unity
does not identify their packets. QA08 and QA09 supply an explicit
rotation and Fourier calculation where a bridge is needed. Neither
statement is claimed as a generic theorem in the inherited Lean package
used here.

\hypertarget{the-exact-modulus-twelve-arithmetic-calculation}{%
\section{The exact modulus-twelve arithmetic
calculation}\label{the-exact-modulus-twelve-arithmetic-calculation}}

We now fix \(q=12\). The coprime labels are

\[
U_{12}=\{1,5,7,11\},
\qquad \varphi(12)=4,
\qquad \zeta=e^{\pi i/6}=\frac{\sqrt3}{2}+\frac i2.
\]

The signal is

\[
c(n)=\zeta^n+\zeta^{5n}+\zeta^{7n}+\zeta^{11n}
=2\cos(\pi n/6)+2\cos(5\pi n/6).
\]

The cosine formula comes from conjugate pairs. It exposes interference,
while the integer divisor formula provides an independent arithmetic
route to each value.

\hypertarget{theorem-qa10.-the-exact-phase-and-transform-tables}{%
\subsection*{Theorem QA10. The exact phase and transform
tables}\label{theorem-qa10.-the-exact-phase-and-transform-tables}}

Over one period,

\[
c=(4,0,2,0,-2,0,-4,0,-2,0,2,0),
\]

and

\[
\widehat c=(0,12,0,0,0,12,0,12,0,0,0,12).
\]

\textbf{Proof.} For \(q=12\), QA04 becomes

\[
c(n)=2\mathbf1_{2\mid n}-4\mathbf1_{4\mid n}
-6\mathbf1_{6\mid n}+12\mathbf1_{12\mid n}.
\]

Indeed, the only divisors \(d\) of twelve for which \(\mu(12/d)\ne0\)
are \(2,4,6,12\), with coefficients \(2,-4,-6,12\). Evaluating the
divisibility indicators at \(0,\ldots,11\) gives the sample table. QA05
gives value twelve at the four coprime labels and zero elsewhere.
\(\square\)

At sample zero, all four phases agree and the value is four. At sample
one, they cancel. At sample six, all four are minus one and the value is
minus four. Those signs are the relative phase information that will
later distinguish two states with the same initial probabilities.

\hypertarget{a-compact-projector-matrix}{%
\subsection{A compact projector
matrix}\label{a-compact-projector-matrix}}

Put \(P=P_{12}^{(12)}\). Its entries are

\[
P_{nm}=\frac{c(n-m)}{12}.
\]

Equivalently, its first row is

\[
\left(\frac13,0,\frac16,0,-\frac16,0,-\frac13,0,-\frac16,0,\frac16,0\right),
\]

and each subsequent row is the cyclic right shift of the preceding row.
This specifies every entry of the twelve-by-twelve rational matrix. It
agrees exactly with the literal matrix in the original formal module and
with the indexed matrix in the later transport module.

With \(D=\operatorname{diag}(\mathbf1_{a\in U_{12}})\) and \(F=F_{12}\),

\[
P=FDF^\dagger,
\qquad P^2=P,
\qquad P^\dagger=P,
\qquad \operatorname{rank}P=\operatorname{tr}P=4.
\]

The six divisor shells of length twelve have orders \(1,2,3,4,6,12\) and
ranks \(1,1,2,2,2,4\). Their ranks add to twelve, and their orthogonal
sum is the identity by QA06.

\hypertarget{the-finite-lanes-of-the-same-signal}{%
\subsection{The finite lanes of the same
signal}\label{the-finite-lanes-of-the-same-signal}}

For \(L=3\) and \(M=4\), the coefficient-residue lanes are

\[
\begin{aligned}
c_0&=(4,0,-4,0),\\
c_1&=(0,-2,0,2),\\
c_2&=(2,0,-2,0).
\end{aligned}
\]

Their four-point Fourier transforms are

\[
\widehat c_0=(0,8,0,8),
\quad \widehat c_1=(0,4i,0,-4i),
\quad \widehat c_2=(0,4,0,4).
\]

For example, a full coefficient at \(k=1\) is reconstructed as

\[
\widehat c(1)=8+4i\zeta^{-1}+4\zeta^{-2}=12.
\]

At \(k=9\), the same lane-frequency index \(1\) gives

\[
\widehat c(9)=8+4i\zeta^{-9}+4\zeta^{-18}=8-4-4=0.
\]

The offsets distinguish two full Fourier labels with the same reduced
frequency. These equations demonstrate the role of the phases in QA09.
They do not identify a lane with a character-order shell.

\hypertarget{filtering-a-signal-with-something-to-remove}{%
\section{Filtering a signal with something to
remove}\label{filtering-a-signal-with-something-to-remove}}

Take the off-band perturbation

\[
b(n)=3\cos(\pi n/2)+2.
\]

The constant occupies frequency zero, and the cosine occupies labels
three and nine. These labels all lie outside \(U_{12}\). The
perturbation has the literal rational vector

\[
b=(5,2,-1,2,5,2,-1,2,5,2,-1,2).
\]

Set \(x=c+b\), giving

\[
x=(9,2,1,2,3,2,-5,2,3,2,1,2).
\]

\Needspace{13\baselineskip}
\hypertarget{theorem-qa11.-exact-filter-recovery-and-orthogonal-energy}{%
\subsection*{Theorem QA11. Exact filter recovery and orthogonal
energy}\label{theorem-qa11.-exact-filter-recovery-and-orthogonal-energy}}

The test signal satisfies

\[
Pc=c,
\qquad Pb=0,
\qquad Px=c,
\qquad (I-P)x=b,
\]

and

\[
\langle c,b\rangle=0,
\qquad \|c\|^2=48,
\qquad \|b\|^2=102,
\qquad \|x\|^2=150=48+102.
\]

\textbf{Proof.} QA05 gives the support of \(\widehat c\). Write the
cosine as

\[
3\cos(\pi n/2)=\frac32(\zeta^{3n}+\zeta^{9n}).
\]

Thus \(\widehat b\) equals twenty-four at zero, eighteen at three and
nine, and zero elsewhere. Multiplication by \(D\) fixes the spectrum of
\(c\) and kills the spectrum of \(b\), proving the filter identities.
Their supports are disjoint, so Parseval for inner products gives
\(\langle c,b\rangle=0\). QA05 gives \(\|c\|^2=12\cdot4=48\), and

\[
\|b\|^2=\frac{24^2+18^2+18^2}{12}=102.
\]

QA07 then gives \(\|x\|^2=48+102=150\). The same numbers follow by
direct rational squaring of the displayed vectors. \(\square\)

The discarded component is part of a complete explanation of the input.
It is called a perturbation because the example declares \(c\) to be the
target. It is a perfectly exact signal component rather than an error
introduced by the representation. If the task were to recover all of
\(x\), deleting \(b\) would lose valid information.

\hypertarget{normalized-states-and-literal-unitary-transport}{%
\section{Normalized states and literal unitary
transport}\label{normalized-states-and-literal-unitary-transport}}

In a twelve-dimensional complex Hilbert space, let
\(|0\rangle,\ldots,|11\rangle\) be an orthonormal basis. Define

\[
|u\rangle=\frac{|1\rangle+|5\rangle+|7\rangle+|11\rangle}{2}.
\]

Its squared norm is \(4(1/2)^2=1\). Define the amplitude states

\[
|\psi\rangle=\frac1{\sqrt{48}}\sum_{n=0}^{11}c(n)|n\rangle,
\qquad
|\chi\rangle=\frac1{\sqrt{48}}\sum_{n=0}^{11}|c(n)|\,|n\rangle.
\]

Both have norm one because \(\sum_nc(n)^2=48\). The second state is a
separately specified state, constructed from the absolute-value vector.
No universal quantum operation implementing the nonlinear rule
\(z\mapsto|z|\) is being asserted.

\Needspace{16\baselineskip}
\hypertarget{theorem-qa12.-exact-preparation-and-recovery}{%
\subsection*{Theorem QA12. Exact preparation and
recovery}\label{theorem-qa12.-exact-preparation-and-recovery}}

With the positive-sign Fourier matrix defined in Section 2,

\[
\boxed{F|u\rangle=|\psi\rangle,
\qquad F^\dagger|\psi\rangle=|u\rangle.}
\]

\textbf{Proof.} At coordinate \(n\), the first left-hand side is

\[
\frac1{2\sqrt{12}}\sum_{a\in U_{12}}\zeta^{na}
=\frac{c(n)}{\sqrt{48}}.
\]

This is the defining coordinate of \(\psi\). Applying \(F^\dagger\) and
using \(F^\dagger F=I\) proves the second equality. \(\square\)

The normalization relation is \(\sqrt{48}=2\sqrt{12}\). In the exact
coordinate implementation the Fourier scale is \(\sqrt3/6=1/\sqrt{12}\),
and the state scale is \(\sqrt3/12=1/\sqrt{48}\). The later formal
source proves those scale identities and the exponential entry formula,
rather than leaving them as informal interpretations of a rational
tuple.

This is a mathematical quantum-state construction. An actual realization
would additionally require state preparation, a twelve-dimensional
unitary implementation, coherence control, and sampling. A
twelve-dimensional Fourier matrix is not the ordinary
sixteen-dimensional Fourier transformation on four qubits. The unitary
direct sum \(F_{12}\oplus I_4\) supplies a mathematical
sixteen-dimensional embedding, but it supplies no gate synthesis or
resource bound by itself.

\hypertarget{relative-phase-is-absent-from-one-histogram-and-present-in-another}{%
\section{Relative phase is absent from one histogram and present in
another}\label{relative-phase-is-absent-from-one-histogram-and-present-in-another}}

For a normalized state \(|v\rangle=\sum_nv(n)|n\rangle\), the
standard-basis probabilities are \(|v(n)|^2\). In the basis analyzed by
\(F^\dagger\), the probabilities are \(|(F^\dagger v)(k)|^2\).

These describe two measurement procedures on separately prepared copies.
Measuring the first histogram on a particular copy does not preserve
that copy as a coherent input for the second procedure.

\hypertarget{theorem-qa13.-equal-initial-probabilities-and-orthogonal-states}{%
\subsection*{Theorem QA13. Equal initial probabilities and orthogonal
states}\label{theorem-qa13.-equal-initial-probabilities-and-orthogonal-states}}

The states \(\psi\) and \(\chi\) have identical probabilities in the
original basis and are orthogonal:

\[
|\psi(n)|^2=|\chi(n)|^2\quad\text{for all }n,
\qquad \langle\psi,\chi\rangle=0.
\]

\textbf{Proof.} Squaring a real number or its absolute value gives the
same result, proving the first claim. For orthogonality,

\[
\langle\psi,\chi\rangle
=\frac1{48}\sum_{n=0}^{11}c(n)|c(n)|.
\]

The positive entries contribute \(16+4+4=24\), and the negative entries
contribute \(-4-16-4=-24\). Their sum is zero. \(\square\)

This is a strong information boundary. The states are distinct even
after the physically irrelevant overall phase is disregarded. Orthogonal
normalized states cannot be scalar multiples of one another. Yet their
original-basis probability histograms are identical.

\Needspace{18\baselineskip}
\hypertarget{theorem-qa14.-complete-arithmetic-phase-separation}{%
\subsection*{Theorem QA14. Complete arithmetic phase
separation}\label{theorem-qa14.-complete-arithmetic-phase-separation}}

After applying \(F^\dagger\),

\[
\Pr_{F^\dagger\psi}(k\in U_{12})=1,
\qquad
\Pr_{F^\dagger\chi}(k\in U_{12})=0.
\]

More explicitly,

\[
(F^\dagger\psi)(k)
=\begin{cases}1/2,&k\in\{1,5,7,11\},\\0,&\text{otherwise},\end{cases}
\]

whereas

\[
(F^\dagger\chi)(k)
=\begin{cases}
2/3,&k\in\{0,6\},\\
1/6,&k\in\{2,4,8,10\},\\
0,&\text{otherwise}.
\end{cases}
\]

\textbf{Proof.} QA12 gives the first amplitude table. For the second,
set \(a(n)=|c(n)|\). Its vector is

\[
a=(4,0,2,0,2,0,4,0,2,0,2,0).
\]

It has period six. Pairing \(n\) and \(n+6\) in its twelve-point
transform produces a factor \(1+(-1)^k\), so every odd transform
coefficient vanishes. For \(k=2\ell\), its transform is twice the
six-point transform of \((4,0,2,0,2,0)\). The doubled transform is
sixteen when \(\ell\) is divisible by three and four otherwise, as one
sees from

\[
2\left(4+2e^{-2\pi i\ell/3}+2e^{-4\pi i\ell/3}\right).
\]

When \(\ell\equiv0\pmod3\), the bracket is eight. Otherwise the two
cubic phases sum to minus one and the bracket is two. Thus

\[
\widehat a=(16,0,4,0,4,0,16,0,4,0,4,0).
\]

State normalization and the unitary inverse transform divide this table
by \(\sqrt{48}\sqrt{12}=24\), producing the second amplitude table. Its
support is even, while every label in \(U_{12}\) is odd. The selected
probability is therefore zero. The signed state has four probabilities
of \(1/4\), with total one. \(\square\)

\Needspace{18\baselineskip}
The full probability comparison is:

{\small\begin{longtable}[]{@{}
  >{\centering\arraybackslash}p{\dimexpr 0.2500\columnwidth - 1.5000\tabcolsep\relax}
  >{\centering\arraybackslash}p{\dimexpr 0.2500\columnwidth - 1.5000\tabcolsep\relax}
  >{\centering\arraybackslash}p{\dimexpr 0.2500\columnwidth - 1.5000\tabcolsep\relax}
  >{\centering\arraybackslash}p{\dimexpr 0.2500\columnwidth - 1.5000\tabcolsep\relax}@{}}
\toprule\noalign{}
\begin{minipage}[b]{\linewidth}\centering
Label
\end{minipage} & \begin{minipage}[b]{\linewidth}\centering
Initial probability, both states
\end{minipage} & \begin{minipage}[b]{\linewidth}\centering
After \(F^\dagger\), signed state
\end{minipage} & \begin{minipage}[b]{\linewidth}\centering
After \(F^\dagger\), unsigned state
\end{minipage} \\
\midrule\noalign{}
\endhead
\bottomrule\noalign{}
\endlastfoot
0 & \(1/3\) & \(0\) & \(4/9\) \\
1 & \(0\) & \(1/4\) & \(0\) \\
2 & \(1/12\) & \(0\) & \(1/36\) \\
3 & \(0\) & \(0\) & \(0\) \\
4 & \(1/12\) & \(0\) & \(1/36\) \\
5 & \(0\) & \(1/4\) & \(0\) \\
6 & \(1/3\) & \(0\) & \(4/9\) \\
7 & \(0\) & \(1/4\) & \(0\) \\
8 & \(1/12\) & \(0\) & \(1/36\) \\
9 & \(0\) & \(0\) & \(0\) \\
10 & \(1/12\) & \(0\) & \(1/36\) \\
11 & \(0\) & \(1/4\) & \(0\) \\
\end{longtable}}

Every column sums to one. The Fourier probability supports are disjoint.
Negative amplitudes in the original arithmetic signal therefore encode
observable relative phase, even though their signs do not appear in the
first probability column.

\hypertarget{theorem-qa15.-no-inverse-from-the-initial-histogram-alone}{%
\subsection*{Theorem QA15. No inverse from the initial histogram
alone}\label{theorem-qa15.-no-inverse-from-the-initial-histogram-alone}}

There is no function from original-basis probability histograms to state
vectors that reconstructs both \(\psi\) and \(\chi\). There is also no
reconstruction of both states up to overall phase from that histogram
alone.

\textbf{Proof.} The two histograms are equal by QA13. A function assigns
the same output to equal inputs, so it cannot output the two distinct
vectors. For the phase-insensitive formulation, QA13 shows that the
states are orthogonal and normalized. If \(\chi=\lambda\psi\) for any
complex scalar of modulus one, then
\(\langle\psi,\chi\rangle=\lambda\ne0\), a contradiction. \(\square\)

The first formulation is formalized in the later publication module. The
strengthening to equivalence up to global phase is a prose corollary
proved here. It is not counted as an additional inherited Lean
declaration.

\hypertarget{the-selector-as-a-projective-measurement}{%
\section{The selector as a projective
measurement}\label{the-selector-as-a-projective-measurement}}

The pair \(\{P,I-P\}\) is a projective measurement because its two
members are self-adjoint, idempotent, orthogonal, and sum to the
identity. For the conditional update we choose the ideal Lüders
instrument, whose unnormalized branches send a density matrix \(\rho\)
to \(P\rho P\) and \((I-P)\rho(I-P)\). The projector effects determine
the outcome probabilities; this chosen instrument determines the state
update. The probability rule and this conditional update can be stated
independently of the arithmetic example.

\hypertarget{theorem-qa16.-selection-probability-and-conditional-state}{%
\subsection*{Theorem QA16. Selection probability and conditional
state}\label{theorem-qa16.-selection-probability-and-conditional-state}}

Let \(P\) be an orthogonal projector on \(\mathbb C^N\) and let \(v\) be
a normalized state. For the ideal Lüders instrument just specified, the
two probabilities are

\[
p=\|Pv\|^2=\langle v,Pv\rangle,
\qquad 1-p=\|(I-P)v\|^2.
\]

When \(p>0\), the normalized selected state is \(Pv/\sqrt p\). When
\(1-p>0\), the normalized complementary state is \((I-P)v/\sqrt{1-p}\).
The probabilities belong to \([0,1]\) and sum to one.

\textbf{Proof.} Since \(P^\dagger P=P\), the quadratic form equals
\(\|Pv\|^2\) and is nonnegative. QA07 applied to a unit vector gives the
probability sum and the upper bound. Division by the positive square
root produces a state of squared norm one in each nonzero branch.
\(\square\)

For the explicit normalized test signal \(v=x/\sqrt{150}\), QA11 gives

\[
p=\frac{48}{150}=\frac8{25},
\qquad 1-p=\frac{102}{150}=\frac{17}{25}.
\]

The selected conditional state is \(c/\sqrt{48}=\psi\), and the
complementary conditional state is \(b/\sqrt{102}\). An exact classical
filter applied to the full vector returns \(c\). A projective
measurement on the normalized input instead returns a branch with these
probabilities. These are compatible uses of the same operator with
different observation rules.

The vector \(\chi\) makes noninjectivity concrete: \(P\chi=0\) although
\(\|\chi\|=1\). A reversible Fourier change of basis cannot be
substituted for this loss of information. Conversely, retaining both
complementary packets reconstructs every vector, as QA07 and the
publication module explicitly state.

The inherited formal source certifies the concrete projectors, states,
and probability identities. It does not provide a general quantum
measurement library or a generic theorem implementing the conditional
state-update rule. QA16 is a general prose theorem, with the usual
finite-dimensional measurement interpretation.

\hypertarget{exact-arithmetic-coordinates-and-their-complex-meaning}{%
\section{Exact arithmetic coordinates and their complex
meaning}\label{exact-arithmetic-coordinates-and-their-complex-meaning}}

The twelfth root contains \(\sqrt3\) and \(i\). Its relevant exact
arithmetic can be represented without floating-point trigonometric
evaluation. Let a tuple

\[
(a,b,c,d)\in\mathbb Q^4
\]

denote

\[
\iota(a,b,c,d)=a+b\sqrt3+i(c+d\sqrt3).
\]

The coordinate root is \((0,1/2,1/2,0)\), whose interpretation is
\(\sqrt3/2+i/2\). Coordinate conjugation sends \((a,b,c,d)\) to
\((a,b,-c,-d)\). Addition is entrywise, and multiplication is

\[
\begin{aligned}
(a,b,c,d)(A,B,C,D)=\big(&aA+3bB-cC-3dD,\\
&aB+bA-cD-dC,\\
&aC+3bD+cA+3dB,\\
&aD+bC+cB+dA\big).
\end{aligned}
\]

\hypertarget{theorem-qa17.-sound-interpretation-of-the-exact-operations}{%
\subsection*{Theorem QA17. Sound interpretation of the exact
operations}\label{theorem-qa17.-sound-interpretation-of-the-exact-operations}}

The map \(\iota\) preserves zero, one, rational scalars, addition,
multiplication, conjugation, natural powers, and finite sums. Therefore
an equality obtained with these exact coordinate operations yields the
corresponding equality in \(\mathbb C\).

\textbf{Proof.} The claims for zero, one, rational scalars, addition,
and conjugation follow directly from the definition. For multiplication,
write the two interpreted values as \((a+b\sqrt3)+i(c+d\sqrt3)\) and
\((A+B\sqrt3)+i(C+D\sqrt3)\). Their product has real part

\[
(a+b\sqrt3)(A+B\sqrt3)-(c+d\sqrt3)(C+D\sqrt3)
\]

and imaginary part

\[
(a+b\sqrt3)(C+D\sqrt3)+(c+d\sqrt3)(A+B\sqrt3).
\]

Expanding and using \((\sqrt3)^2=3\) gives exactly the displayed tuple
product. Power preservation follows by induction from multiplication and
one. Finite-sum preservation follows by induction from addition and
zero. Applying a function to equal coordinate tuples gives equal complex
numbers. \(\square\)

Only sound transport of equalities is needed here. The formal source
does not install the tuple type as a global field instance, and it does
not need an injectivity theorem for \(\iota\). It proves preservation of
the actual operations and explicitly connects the scales and root to
their conventional complex formulas.

The matrix interpretation is entrywise. Since multiplication is a finite
sum of products and the adjoint uses conjugation, QA17 also transports
coordinate matrix multiplication and adjoints. The same argument
transports matrix action on vectors and squared norms. This is what
closes the gap between an exact coordinate certificate and an identity
about an ordinary complex unitary matrix.

\hypertarget{the-central-finite-bridge-theorem}{%
\section{The central finite bridge
theorem}\label{the-central-finite-bridge-theorem}}

The arithmetic, signal, and state statements can now be gathered without
leaving any undefined arrow.

\hypertarget{theorem-qa18.-arithmetic-rational-signal-and-unitary-realization}{%
\subsection*{Theorem QA18. Arithmetic, rational signal, and unitary
realization}\label{theorem-qa18.-arithmetic-rational-signal-and-unitary-realization}}

Let \(F\), \(D\), \(P\), \(c\), \(b\), \(x\), \(u\), \(\psi\), and
\(\chi\) have the explicit modulus-twelve definitions above. Then:

\begin{enumerate}
\def\labelenumi{\arabic{enumi}.}

\item
  \(F^\dagger F=FF^\dagger=I\).
\item
  \(P=FDF^\dagger\), \(P^2=P\), and \(P^\dagger=P\).
\item
  \(Px=c\) and \((I-P)x=b\), with \(\|x\|^2=48+102=150\).
\item
  \(Fu=\psi\) and \(F^\dagger\psi=u\).
\item
  \(u\), \(\psi\), and \(\chi\) are normalized, and \(\psi\) is
  orthogonal to \(\chi\).
\item
  \(\psi\) and \(\chi\) have identical original-basis probabilities.
\item
  Their inverse-Fourier coprime probabilities are respectively one and
  zero.
\item
  \(P\psi=\psi\), \(P\chi=0\), and \(P\) is not injective.
\end{enumerate}

\textbf{Proof.} QA03 gives the unitarity identities. QA06 gives
projector transport and the two projector identities. QA11 gives the
filter and energy statements. QA12 gives preparation and recovery; the
displayed normalizations and QA13 give statement five and statement six.
QA14 gives statement seven. In Fourier coordinates \(\psi\) is supported
on \(U_{12}\) while \(\chi\) is supported outside it, so \(P\psi=\psi\)
and \(P\chi=0\). The normalized nonzero vector \(\chi\) is therefore in
the kernel, proving noninjectivity. \(\square\)

The later formal publication theorem gathers the principal complex
transport and probability identities into a single conjunction. Its
supporting modules also supply the exact vector coefficients, filter
recovery, residual reconstruction, and noninjectivity theorem. The
general prose route above explains why the identities hold; the finite
formal source certifies this literal example under its exact
definitions.

One mathematical object has thus appeared in three realizations.
Arithmetic selects coprime character labels. The signal realization is a
rational circulant projector. The complex state realization uses a
unitary basis change and a projective measurement. The complete vector
information is preserved by the basis change, while selection and
observation have separately stated losses.

\hypertarget{a-continuous-phase-family-with-a-fixed-initial-histogram}{%
\section{A continuous phase family with a fixed initial
histogram}\label{a-continuous-phase-family-with-a-fixed-initial-histogram}}

The two endpoint states are especially simple, but the mechanism is
continuous. Define the nonnegative vectors \(v_+\) and \(v_-\) by
placing \(|c(n)|/\sqrt{48}\) respectively on the positive and negative
entries of \(c\), with zero elsewhere. Thus

\[
\chi=v_++v_-,
\qquad \psi=v_+-v_-.
\]

Their supports are disjoint, and each has squared norm \(1/2\). For real
\(\theta\), define

\[
\chi_\theta=v_++e^{i\theta}v_-.
\]

It is a normalized state with the same original-basis probabilities for
every \(\theta\). The endpoints are \(\chi_0=\chi\) and
\(\chi_\pi=\psi\).

\hypertarget{theorem-qa19.-relative-phase-response}{%
\subsection*{Theorem QA19. Relative-phase
response}\label{theorem-qa19.-relative-phase-response}}

For this family,

\[
\boxed{\Pr_{F^\dagger\chi_\theta}(k\in U_{12})
=\|P\chi_\theta\|^2=\sin^2(\theta/2).}
\]

\textbf{Proof.} Disjoint support gives
\(\|\chi_\theta\|^2=\|v_+\|^2+\|v_-\|^2=1\) and unchanged entrywise
squared magnitudes. Solving for \(v_+\) and \(v_-\) in terms of \(\chi\)
and \(\psi\) gives

\[
\chi_\theta
=\frac{1+e^{i\theta}}2\chi
+\frac{1-e^{i\theta}}2\psi.
\]

Since \(P\chi=0\) and \(P\psi=\psi\),

\[
P\chi_\theta=\frac{1-e^{i\theta}}2\psi.
\]

Its squared norm is

\[
\frac{|1-e^{i\theta}|^2}{4}
=\frac{1-\cos\theta}{2}
=\sin^2(\theta/2).
\]

Finally, \(P=FDF^\dagger\) and unitarity imply
\(\|P\chi_\theta\|^2=\|DF^\dagger\chi_\theta\|^2\), exactly the sum of
the selected Fourier-basis probabilities. \(\square\)

This is a proved prose extension for the visual guide. It is not
presented as an inherited Lean theorem. It shows that a probability
histogram can stay fixed while an interference-based observation varies
from zero to one.

\hypertarget{phase-replacement-and-dephasing-are-different-operations}{%
\section{Phase replacement and dephasing are different
operations}\label{phase-replacement-and-dephasing-are-different-operations}}

The unsigned endpoint is a coherent state with specified positive
amplitudes. It should not be confused with a state obtained by
forgetting all relative phases. The latter is represented by a mixed
state, not generally by one state vector.

For a normalized vector \(v\), its density matrix is \(vv^\dagger\).
Complete dephasing in the sample basis keeps only its diagonal entries.
In the present example, dephasing either \(\psi\) or \(\chi\) gives the
same diagonal matrix

\[
\rho_{\mathrm{diag}}=\operatorname{diag}\big(c(n)^2/48\big)_{n=0}^{11}.
\]

Its trace is one, and every diagonal entry is nonnegative.

\hypertarget{theorem-qa20.-the-dephased-selected-probability}{%
\subsection*{Theorem QA20. The dephased selected
probability}\label{theorem-qa20.-the-dephased-selected-probability}}

The selected projective-measurement probability of this diagonal state
is

\[
\operatorname{tr}(P\rho_{\mathrm{diag}})=\frac13.
\]

\textbf{Proof.} Every diagonal entry of \(P\) is \(c(0)/12=4/12=1/3\).
Multiplication by a diagonal matrix and taking the trace gives

\[
\operatorname{tr}(P\rho_{\mathrm{diag}})
=\sum_n P_{nn}\frac{c(n)^2}{48}
=\frac13\sum_n\frac{c(n)^2}{48}
=\frac13.
\] \(\square\)

The three outcomes are therefore distinct: the signed coherent state has
selected probability one, the unsigned coherent state has probability
zero, and the completely dephased state has probability one third. This
calculation prevents an incorrect identification of amplitude sign
removal with physical decoherence. QA20 is a general finite-dimensional
calculation specialized to the example and proved here. Density matrices
and dephasing are outside the inherited Lean package.

\hypertarget{formal-coverage-lineage-and-reproducibility}{%
\section{Formal coverage, lineage, and
reproducibility}\label{formal-coverage-lineage-and-reproducibility}}

\hypertarget{two-inherited-checkpoints-and-the-expanded-boundary}{%
\subsection{Two inherited checkpoints and the expanded
boundary}\label{two-inherited-checkpoints-and-the-expanded-boundary}}

The earlier checked package has four mathematical production modules and
fifty-nine named theorem declarations. Its first module preserves
twenty-two original statements. Its proof-boundary document correctly
says that explicit matrix transport between independently defined
representations and normalized state-action identities are not supplied
as additional theorems in that package. Those limitations must not be
erased when describing the earlier evidence.

\Needspace{16\baselineskip}
The later checked package has five production modules and ninety-six
named theorem declarations:

{\small\begin{longtable}[]{@{}
  >{\raggedright\arraybackslash}p{\dimexpr 0.3333\columnwidth - 1.3332\tabcolsep\relax}
  >{\raggedleft\arraybackslash}p{\dimexpr 0.3333\columnwidth - 1.3332\tabcolsep\relax}
  >{\raggedright\arraybackslash}p{\dimexpr 0.3333\columnwidth - 1.3332\tabcolsep\relax}@{}}
\toprule\noalign{}
\begin{minipage}[b]{\linewidth}\raggedright
Module
\end{minipage} & \begin{minipage}[b]{\linewidth}\raggedleft
Named theorems
\end{minipage} & \begin{minipage}[b]{\linewidth}\raggedright
Principal role
\end{minipage} \\
\midrule\noalign{}
\endhead
\bottomrule\noalign{}
\endlastfoot
\texttt{QuantumArithmeticBridge} & 22 & Literal rational projector and
cyclotomic tables \\
\texttt{FiniteFourier} & 17 & Indexed exact Fourier and normalized state
identities \\
\texttt{SemanticBridge} & 30 & Sound complex interpretation and
transported identities \\
\texttt{Conventions} & 17 & Exponential entries, scales, and complex
projector laws \\
\texttt{Publication} & 10 & Phase lookup, residual reconstruction,
nonreconstruction, and central bundle \\
Total & 96 & Finite modulus-twelve proof chain, including helpers \\
\end{longtable}}

The later source explicitly supplies \(FDF^\dagger=P\), both normalized
state-action identities, the conventional exponential formula for \(F\),
identical initial probabilities, orthogonality, exact selected
probability one versus zero, histogram nonreconstruction, and projector
noninjectivity. It therefore closes the representation-transport
obligations that the earlier package leaves outside its formal boundary.

The ninety-six count includes helper lemmas and bundled propositions. It
is not a count of ninety-six independent discoveries. Nor should the
earlier fifty-nine and later ninety-six counts be added as if their
shared declarations were disjoint.

\hypertarget{current-and-inherited-checks}{%
\subsection{Current and inherited
checks}\label{current-and-inherited-checks}}

This edition preserves the source, theorem index, verification records,
and compiled-artifact lineage of the checked checkpoints. The archived
records provide the earlier compiler and kernel-checking evidence.
Current static audits and independent Python replay provide current
bounded evidence about the downloaded files and the example
implementation. These are disclosed separately in the release records.

A fresh Lean runtime could not be restored in this edition. The
attempted runtime retrieval failed with server errors on three retries,
and no Lean compiler was run. Accordingly, this paper does not describe
the archived Lean evidence as a new build. A future fresh build must use
the checkpoint's exact pinned toolchain and dependencies and record its
own result.

The exact source uses kernel-reduction decision proofs for finite
identities and ordinary proof terms for the transport lemmas. The
declaration audit rejects admitted proofs and unapproved
project-specific axioms. Standard foundational dependencies are reported
in the formal guide. Compiled artifact checking reuses Lean's kernel; it
is not an independently implemented proof assistant.

The later checkpoint preserves three negative controls with distinct
meanings. The false-sign claim is a directly false finite statement. The
false-mode case is an operational rejection whose failure mechanism must
be read alongside its diagnostic and the positive checked transform
table. The admitted-proof case tests the verification policy for
incomplete proof evidence. These controls are not combined with the four
controls of the earlier checkpoint into an artificial total.

\hypertarget{independent-arithmetic-evidence}{%
\subsection{Independent arithmetic
evidence}\label{independent-arithmetic-evidence}}

The original reproducibility package records 3,003 exact Python
assertions, separately from floating-point checks. Its exact verifier
represents the twelfth root in a cyclotomic quotient, checks the signal
and transform tables, the rational projector identities, the test-signal
decomposition, divisor-shell completeness, and the phase-sensitive
example. Its arithmetic regression range is bounded to moduli two
through sixty-four.

The release also includes the current replay and any added finite checks
with their own records. Exact rational or cyclotomic arithmetic and
numerical residuals are reported separately. A small floating-point
residual verifies the implementation at its declared tolerance; it does
not replace the exact proof of the general formula. Likewise, a finite
modulus sweep does not prove the arbitrary-modulus theorem. QA01 through
QA09 provide the general proofs needed here.

\hypertarget{the-twenty-two-preserved-original-statements}{%
\subsection{The twenty-two preserved original
statements}\label{the-twenty-two-preserved-original-statements}}

The original bridge module's statements remain available under their
original names. The following list records their substantive content
rather than treating each source helper as a separate main theorem.

{\small\begin{longtable}[]{@{}
  >{\raggedright\arraybackslash}p{\dimexpr 0.5000\columnwidth - 1.0000\tabcolsep\relax}
  >{\raggedright\arraybackslash}p{\dimexpr 0.5000\columnwidth - 1.0000\tabcolsep\relax}@{}}
\toprule\noalign{}
\begin{minipage}[b]{\linewidth}\raggedright
Original declaration
\end{minipage} & \begin{minipage}[b]{\linewidth}\raggedright
Preserved statement
\end{minipage} \\
\midrule\noalign{}
\endhead
\bottomrule\noalign{}
\endlastfoot
\texttt{ramanujan12\_length} & The rational signal list has length
twelve \\
\texttt{projector12\_shape} & The literal matrix has twelve rows and
twelve columns \\
\texttt{projector12\_symmetric} & The rational matrix equals its
transpose \\
\texttt{projector12\_idempotent} & Exact rational multiplication gives
\(P^2=P\) \\
\texttt{projector12\_fixes\_ramanujan} & The matrix fixes the signal
\(c\) \\
\texttt{projector12\_kills\_noise} & The matrix sends the displayed
perturbation \(b\) to zero \\
\texttt{projector12\_recovers\_test\_signal} & Applying it to \(x\)
returns \(c\) \\
\texttt{ramanujan12\_energy} & \(\|c\|^2=48\) \\
\texttt{noise12\_energy} & \(\|b\|^2=102\) \\
\texttt{testSignal12\_energy} & \(\|x\|^2=150\) \\
\texttt{energy\_decomposition} & The exact energy balance is
\(150=48+102\) \\
\texttt{zeta12\_order\_certificate} & The coordinate root raised to the
twelfth power is one \\
\texttt{zeta12\_inverse\_certificate} & The coordinate root times its
conjugate is one \\
\texttt{ramanujanC12\_table} & Four primitive phase sums give the exact
twelve-entry signal table \\
\texttt{ramanujan\_dft\_table} & The negative-sign DFT gives twelve
exactly at \(1,5,7,11\) \\
\texttt{abs\_ramanujan\_dft\_table} & The unsigned DFT gives the exact
even-support table \\
\texttt{same\_computational\_basis\_probabilities} & Signed and unsigned
rational entries have the same squares \\
\texttt{abs\_ramanujan12\_energy} & The unsigned signal also has squared
energy forty-eight \\
\texttt{recovered\_mode\_probability} & The scalar normalization gives
\(12^2/(48\cdot12)=1/4\) \\
\texttt{recovered\_total\_probability} & Four recovered mode
probabilities sum to one \\
\texttt{phase\_stripping\_coprime\_probability} & Zero transformed
coefficients give zero selected mass \\
\texttt{arithmetic\_signal\_quantum\_bridge} & The original finite
projector, recovery, DFT, and equal-square claims are bundled \\
\end{longtable}}

The later source strengthens the specification around these statements.
It proves that the four labels are exactly the coprime residues, that
the coordinate root has no smaller positive order, that the literal
rational and indexed matrices agree, and that the probability
consequences belong to the actual normalized complex state vectors.

The list operations in the original module use explicit finite literal
lists. Their generic helper definitions are not an assertion of
arbitrary-dimensional linear algebra for unsized lists, where zip
operations could truncate. General vector identities in this paper use
fixed finite index sets, and the later matrix source uses
\texttt{Fin\ 12} for exactly that reason.

\hypertarget{what-is-complete-and-what-remains-a-separate-question}{%
\section{What is complete and what remains a separate
question}\label{what-is-complete-and-what-remains-a-separate-question}}

The finite worked bridge is complete as a mathematical result. Its maps
are defined, its identities are proved, its literal coefficients agree
across the rational and cyclotomic representations, and its principal
modulus-twelve transport and phase claims have inherited formal
certificates. The current release provides the traditional paper, visual
guide, formal concordance, source, and reproducibility records.

The general prose theorems also establish Fourier unitarity, the Möbius
divisor formula, rational divisor-projector completeness, and polyphase
aliasing for the declared finite or formal models. Their
arbitrary-modulus formulations are not claimed as generic Lean
declarations in this package. The continuous phase and dephased-state
extensions are proved prose results rather than inherited formal
additions. The visual guide illustrates the continuous phase family.

Several broader questions have different proof obligations. Efficient
preparation of coprime-support states needs an algorithm and a cost
model. Efficient circuit synthesis of the twelve-dimensional
transformation or its qubit embedding needs a gate construction and
resource bounds. Incomplete-sample recovery needs an observation map
that is injective on the declared model, or a stability theorem with
stated conditioning. Noisy implementation needs a perturbation analysis.
A physical claim needs an experimental protocol and measured data.

These questions can be asked using the same operator language. Their
existence does not weaken the exact bridge proved here, and the bridge
does not answer them automatically. The present result shows, in one
fully specified example, how arithmetic selection becomes rational
filtering and phase-sensitive unitary realization while keeping
reconstruction, information loss, and measurement distinct.

\hypertarget{source-lineage}{%
\section{Source lineage}\label{source-lineage}}

The four primary literature entries are provided in the accompanying
reference fragment. Ramanujan's work supplies the arithmetic lineage.
Vaidyanathan's papers establish the signal-processing setting of
Ramanujan spaces and filter banks. The internal Komorebi lineage is
\emph{Stratified Universality for Resolved Representations}, v1.16.0,
Proposition 21.11, printed page 189, with Boundary 21.12 on the
distinction between character orders and coefficient-residue lanes.

The local authoritative source set for the finite example consists of
\texttt{Worked\_Example.md}, the original reproducibility verifier, the
earlier four-module checked source and proof-boundary document, and the
later five-module checked source with its verification records. Their
exact filenames, source hashes, and evidence classification are listed
in the release provenance and the Lean guide. The traditional paper
supplies self-contained mathematical proofs, while those companion
records make the implementation and formal scope inspectable.

% Primary-source attribution checked 2026-10-07 UTC.
% Journal Parts I and II: metadata authenticated; no uninspected equation numbers.
\begin{thebibliography}{99}
\bibitem{Ramanujan1918}
S. Ramanujan,
\emph{On certain trigonometrical sums and their applications in the theory of numbers},
Transactions of the Cambridge Philosophical Society \textbf{22}, no.~13 (1918), 259--276.
Original-paper scan index: \url{https://www.imsc.res.in/~rao/ramanujan/CamUnivCpapers/Cpaper21/page1.htm}.

\bibitem{Vaidyanathan2014FIR}
P. P. Vaidyanathan,
\emph{Ramanujan-sum expansions for finite duration (FIR) sequences},
2014 IEEE International Conference on Acoustics, Speech and Signal Processing (ICASSP), 4933--4937.
\url{https://doi.org/10.1109/ICASSP.2014.6854540}.
Author manuscript: \url{https://systems.caltech.edu/dsp/PPVSquare.pdf}.

\bibitem{Vaidyanathan2014I}
P. P. Vaidyanathan,
\emph{Ramanujan Sums in the Context of Signal Processing---Part I: Fundamentals},
IEEE Transactions on Signal Processing \textbf{62}(16) (2014), 4145--4157.
\url{https://doi.org/10.1109/TSP.2014.2331617}.

\bibitem{Vaidyanathan2014II}
P. P. Vaidyanathan,
\emph{Ramanujan Sums in the Context of Signal Processing---Part II: FIR Representations and Applications},
IEEE Transactions on Signal Processing \textbf{62}(16) (2014), 4158--4172.
\url{https://doi.org/10.1109/TSP.2014.2331624}.
\bibitem{KomorebiStratified16}
Komorebi P.,
\emph{Stratified Universality for Resolved Representations},
v1.16.0, Block 16, internal manuscript, 2026.
Proposition 21.11 and Boundary 21.12, printed pp.~189--190.
Preserved source: \texttt{Komorebi\_Stratified\_Universality\_Foundations\_v1.16.0\_Block16.pdf}.
\end{thebibliography}

